\documentclass[11pt]{article}
\usepackage[a4paper,margin=25mm]{geometry}
\usepackage{amsmath,amssymb,amsthm}
\newtheorem{theorem}{Theorem}
\newtheorem{lemma}{Lemma}
\title{Structural certificates for the $S_{10}$ atlas table}
\author{L054 finite application}\date{}
\begin{document}\maketitle

\begin{theorem}
Relative to the acquired numerical catalogue and structural registry,
the $1593$ rows of the current companion table headed $S_{10}$ have the
displayed structural descriptions. An extension expression asserts a normal
subgroup, quotient and the stated splitting behaviour. It need not
specify a unique abstract isomorphism class.
\end{theorem}

\paragraph{Source and acquired results.}
The companion table by Thomas Connor and Dimitri Leemans is dated
27 November 2012 and accompanies \cite{atlas}. Its numerical catalogue,
source-row bijection, actual embedded generators and subgroup orders are
acquired from the separately accepted computation in degree $10$.
The acquired source contains $1593$ conjugacy classes and $8023$ weighted
maximal-subgroup edges. The subgroup census and incidence counts are not
recomputed here; proper-group orders are checked again when constructing
literal transports. The registry contains $676$ finite structural models reconstructed
from accepted earlier applications; the exact files read to reconstruct
it are bound to those reviews. Existing hypotheses and qualifications
remain part of the acquired statements.

The source article's Section 3.2 gives the convention of left association
for an unparenthesized sequence of extensions. This convention is acquired
from its separately reviewed native reading. A new audit joins the native
powers, subscripts, parentheses and operators of every structural cell to
its expression and order. The table has $103$ pages and $38011$ retained
text states. Native pages $1,6,33,52,102,103$ were directly inspected.
This does not amount to visual inspection of all pages. The position audit
uses the same pypdf decoder as the flat extraction, rather than an
independent extraction engine. Rows $77$ and $576$ put redundant parentheses
around the single atom $D_{10}$; only that exact glyph sequence is removed
when comparing their expressions. The raw positions are retained. The
pinned source edition, rather than every historical edition, specifies the
claim.

\begin{lemma}
Let $H$ be generated by specified permutations, let $S$ be the union of
their supports, and put $n=|S|$. Suppose that the order of this same
embedded group is known. If $|H|=n!$, then $H=\operatorname{Sym}(S)$.
If every specified generator is even and $|H|=n!/2$, then
$H=\operatorname{Alt}(S)$.
\end{lemma}
\begin{proof}
Every generator fixes the complement of $S$, so $H$ is a subgroup of
$\operatorname{Sym}(S)$. Equality of finite orders proves the first claim.
The sign map is a homomorphism. If all generators are even, then
$H\leq\operatorname{Alt}(S)$, and equality of orders again gives equality
of groups. No transitivity assumption is needed. In the application the
source row is first joined to its accepted numerical record and exact
generator tuple. The acquired order is then the order of the group in
this lemma; it is not inferred from the printed name.
\end{proof}

This identifies rows
$1,2,3,4,7,11,29,45,109,191,451,738,1397,1547$ directly.
The other $1579$ source occurrences are handled using the acquired
models, $172$ new finite model certificates and complete permutation
transports. The model roster consists
of $87$ split extensions, $69$ direct products and $16$ nonsplit extensions. In particular, an
abstractly symmetric or alternating subgroup with another action is not
accepted by the support lemma unless its actual support and order meet
the stated conditions.

\begin{lemma}
Let $M$ and $H$ be finite groups given by complete closures of specified
permutation generators. A bijection $f:M\to H$ satisfying $f(1)=1$ and
$f(xg)=f(x)f(g)$ for every $x\in M$ and every specified generator $g$
is an isomorphism.
\end{lemma}
\begin{proof}
Every element is a positive word in the generators, because the inverse
of a generator in a finite group is a nonnegative power of it. Induction
on such a word gives $f(xy)=f(x)f(y)$. The bijection is therefore an
isomorphism. The reader constructs proposed images along the complete
closure, checks consistency whenever a word reaches an old element, and
checks every right-generator edge. It compares the entire image with the
actual target subgroup. Merely mapping a generating set is insufficient.
\end{proof}

\paragraph{Products and extensions.}
For a new model $H$, the reader closes the displayed generators of $N$,
checks $N\triangleleft H$, and transports an earlier proved model onto
$N$. A split certificate also supplies $K$, its structure,
$N\cap K=1$ and $|N||K|=|H|$. Normality makes $NK$ a subgroup;
its order then gives $H=NK$. For a colon, the retained convention
requires nontrivial conjugation of $K$ on $N$, witnessed by a pair of
noncommuting generators. A direct product requires every pair of factor
generators to commute, which implies that the factors commute elementwise.
The finite transport and extension predicates are reused from the
accepted common method. Every newly proposed model is proved before it
is used, including models obtained in an earlier candidate search.

Twelve nonsplit models have a prime quotient of order $p$. A complement
exists exactly when there is an element of order $p$ outside $N$: its
cyclic subgroup intersects $N$ trivially, and the order formula gives all
of $H$. The nonsplitting reader checks every outside element, so the
absence of such an element proves the assertion.

The other four nonsplit models have quotient $S_3$. Their specified
quotient generators satisfy $a^2=b^3=(ab)^2=1$ and generate the quotient.
These are the standard defining relations of $S_3$. A complement would
lift them to elements satisfying the same relations. Conversely such
lifts give a homomorphism from $S_3$ whose composite with the quotient
map is the identity, hence a section. All lifts are considered; those
failing the first two relations may be discarded before testing the
third. The recorded count includes the full original Cartesian product.
Each possible pair fails at least one relation, so no section exists.

\paragraph{The proposal search.}
The first pass stopped at row $34$, with expression $A_5\times A_5$:
GAP's general complement routine does not handle the case where both the
normal factor and the quotient are nonsolvable. A process exit code of
zero after the resulting break-loop error is recorded as an incomplete
search. The retained witnesses are not discarded or accepted on that basis.

The following elementary criterion supplies a suitable direct-product
search. Let $N\triangleleft H$ and $n=[H:N]$. If $C_H(N)$ has order $n$
and intersects $N$ trivially, it is the unique direct complement of $N$.
Indeed, a direct complement lies in that centralizer and has the same
order. Otherwise one can inspect the normal subgroups $K$ of order $n$
with $N\cap K=1$. Normality puts every commutator in the intersection,
so these factors commute, and the order formula gives $H=NK$.
Conversely every internal direct complement is normal, since conjugation
by $N$ fixes it pointwise and conjugation by itself preserves it.
This standard direct-product criterion changes proposal generation only.
The literal proof predicates remain unchanged.

\paragraph{Verification and attribution.}
The numerical catalogue, source expressions and association convention
belong to the source and acquired results. The support criterion and the
internal direct-product criterion above are elementary group theory.
The contribution of this application is an explicit source join and
finite certificates for its remaining occurrences using the common
method. The complete literal check passed in $22.115$ seconds, with
$1593$ source joins, $14$ support identifications, $172$ model proofs,
$1907$ transports and $3209070$ generator edges. Nonsplitting requires
$10236$ outside-element tests and $5764$ quotient-lift pairs. Inputs
remained unchanged and stderr was empty. The initial incomplete candidate
run is retained; the $27$ remaining rows were subsequently supplied in
$2.531$ seconds using the direct-factor improvement. Candidate generation
itself is never proof.

Ten damaged source inputs are rejected, and eight narrow glyph checks
confirm that only the two parenthesized $D_{10}$ atoms are normalized.
Support recognition has its separate eight small literal model controls
and nine damaged cases. Ordinary and optimized modes are checked
separately. The final independent review records source admission and
retains the extraction limitations stated above. No conjecture resolution,
Lean formalization, historical-byte identity, whole-article completion or
compression measurement is claimed. Other large companion tables remain
in scope for L054.

\begin{thebibliography}{1}
\bibitem{atlas} T. Connor and D. Leemans,
An atlas of subgroup lattices of finite almost simple groups,
\emph{Ars Mathematica Contemporanea} \textbf{8} (2015), 259--266.
The retained companion table is headed $S_{10}$ and dated
27 November 2012.
\end{thebibliography}
\end{document}
